\chapter{Operational Guide and Reproducibility Manual}
\label{app:operational}

This appendix documents how to assemble the environment, run the pipeline, and
interpret the outputs so that the experiments of Chapter~\ref{ch:experiments}
can be reproduced.

\section{Hardware and System Prerequisites}
\label{sec:app-hardware}

\begin{itemize}
  \item A CUDA-capable GPU for the CLIP and VLM stages (experiments used a
  Kaggle NVIDIA~T4 with 16\,GB memory); a commodity multi-core CPU suffices
  for frame filtering and tracking.
  \item Python~3.10 or later with the package set of
  Section~\ref{sec:app-env}.
  \item Sufficient disk for model weights (YOLOv8-n, CLIP ViT-B/32, and
  Moondream2) and for per-run output logs.
\end{itemize}

\section{Environment Assembly and Dependency Registry}
\label{sec:app-env}

The core dependencies and their roles are listed in
Table~\ref{tab:dependencies}. A minimal setup is shown below.

\begin{lstlisting}[language=bash, caption={Environment setup.}]
python -m venv .venv && source .venv/bin/activate   # or .venv\Scripts\activate on Windows
pip install torch torchvision opencv-python
pip install ultralytics            # YOLOv8
pip install boxmot                 # ByteTrack
pip install git+https://github.com/openai/CLIP.git
pip install transformers           # Moondream2 weights via HF hub
\end{lstlisting}

\begin{table}[htbp]
\centering
\caption[Dependency registry]{Primary software dependencies.}
\label{tab:dependencies}
\small
\begin{tabular}{@{}lll@{}}
\toprule
\textbf{Package} & \textbf{Component} & \textbf{Role} \\
\midrule
\texttt{torch}, \texttt{torchvision} & PyTorch & Tensor / model runtime \\
\texttt{opencv-python}               & OpenCV  & Frame I/O and filtering \\
\texttt{ultralytics}                 & YOLOv8-n & Object detection \\
\texttt{boxmot}                      & ByteTrack & Multi-object tracking \\
\texttt{CLIP}                        & CLIP ViT-B/32 & Semantic embedding \\
\texttt{transformers}                & Moondream2 & Vision-language model \\
\bottomrule
\end{tabular}
\end{table}

\section{Command-Line Execution}
\label{sec:app-cli}

\subsection{Running all baselines and the proposed scheduler}
\label{sec:app-baseline}

The five policies --- \texttt{every\_frame}, \texttt{fixed\_interval}
($K=30$), \texttt{motion\_gating}, \texttt{embedding\_threshold}, and
\texttt{threshold} (the proposed three-tier scheduler) --- are run on an
identical stream by the \texttt{run\_all\_baselines.py} driver. To obtain the
measured metrics used in Chapter~\ref{ch:experiments}, VLM execution must be
\emph{enabled} (omit \texttt{--skip\_vlm}) and the frame budget left
unbounded:

\begin{lstlisting}[language=bash, caption={Run all five policies on the same video (VLM enabled, all frames).}]
python scripts/run_all_baselines.py \
    --video /path/to/video.mp4 \
    --max_frames -1 \
    --device cuda
\end{lstlisting}

\subsection{Individual runs and the threshold sweep}
\label{sec:app-proposed}

\begin{lstlisting}[language=bash, caption={Individual policy runs and the ablation sweep.}]
# Proposed three-tier scheduler (policy "threshold")
python scripts/run_pipeline.py \
    --video /path/to/video.mp4 \
    --policy threshold \
    --tau_low 0.15 --tau_high 0.40 \
    --alpha 0.5 --beta 0.3 --gamma 0.2 \
    --decay_lambda 0.90 \
    --vlm moondream \
    --max_frames -1 \
    --experiment_name proposed_threshold \
    --output_dir experiments/proposed_threshold/

# A single baseline, e.g. embedding-threshold
python scripts/run_pipeline.py \
    --video /path/to/video.mp4 \
    --policy embedding_threshold --tau 0.30 \
    --experiment_name baseline_embed_threshold \
    --output_dir experiments/baseline_embed_threshold/

# Sweep the lower threshold for the ablation
python scripts/sweep_thresholds.py --param tau_low --range 0.05 0.35 0.05
\end{lstlisting}

\section{Output Structure}
\label{sec:app-output}

Each run writes to its \texttt{--out} directory:

\begin{itemize}
  \item \texttt{metrics.json} --- aggregate metrics (VLM call rate, latency,
  throughput, peak memory, retention, SCE).
  \item \texttt{frame\_log.csv} --- per-frame drift components $D_t$,
  $D_{\text{embed}}$, $D_{\text{objects}}$, $D_{\text{track}}$, the chosen
  tier, and latency; the input consumed by the Semantic Observatory.
  \item \texttt{outputs/} --- delivered VLM captions (cached or freshly
  generated) for qualitative inspection.
\end{itemize}

The aggregation script collates every run's \texttt{metrics.json} into the
comparison table. The Semantic Observatory
(Section~\ref{sec:arch-observatory}) is then run post hoc on any
\texttt{frame\_log.csv} to produce the semantic-evolution analytics, figures,
and dashboard:

\begin{lstlisting}[language=bash, caption={Post-hoc Semantic Observatory analysis.}]
python scripts/run_observatory.py \
    --frame_log experiments/proposed_threshold/frame_log.csv \
    --output_dir results/observatory/proposed/ \
    --experiment_name proposed_threshold
\end{lstlisting}
